\section{Introduction}
\label{sec:introduction}

In safety-critical cyber-physical systems, such as commercial aviation turbofan engines, unexpected component failure incurs catastrophic operational, financial, and life-safety hazards. Consequently, industrial fleet operators have historically adhered to rigid preventive maintenance schedules, overhauling major subsystems at fixed operational intervals (flight hours or cycles). While preventive scheduling mitigates catastrophic in-flight failure, it imposes severe economic costs by routinely discarding components with substantial remaining physical life~\cite{saxena2008damage,shehadeh2024evaluating}.

Predictive Maintenance (PdM) leverages onboard multi-sensor telemetry to transition from time-based maintenance to condition-based, opportunistic intervention. Machine learning prognostics aims to estimate the Remaining Useful Life (RUL) of degrading machinery, enabling timely maintenance before failure occurs. In recent years, deep neural architectures—including Convolutional Neural Networks (CNNs), Long Short-Term Memory (LSTM) recurrent networks, and gradient-boosted decision trees—have achieved remarkable point-prediction accuracy on standardized benchmarks such as NASA C-MAPSS.

However, deterministic point predictors $\hat{y}$ provide no quantification of predictive uncertainty. In sequential decision environments, a deterministic forecast is fundamentally incomplete: an operator cannot determine whether a predicted life of $\hat{y} = 15$ cycles carries high confidence or a substantial probability of immediate breakdown at cycle 2. Furthermore, standard models are trained under symmetric loss functions (such as Mean Squared Error), which penalize early and late errors equally. In physical operations, this symmetry is catastrophic: predicting that an engine will survive longer than it actually does ($d > 0$) risks uncontained in-service failure, whereas predicting failure prematurely ($d < 0$) merely incurs early maintenance cost~\cite{chen2022data,diao2026turbofan}.

To provide principled uncertainty bounds, prognostic models must emit calibrated prediction intervals $[L(x), U(x)]$ targeting nominal coverage $1 - \alpha$. Predictive uncertainty methods differ in their assumptions and calibration behavior. Their maintenance value must be evaluated through interval reliability and downstream decision outcomes under a common protocol. Conformal prediction~\cite{javanmardi2023conformal} provides a distribution-free uncertainty framework.

IntelliTwin uses an engine-weighted conformal construction motivated by hierarchical conformal prediction. Formal coverage results require the hierarchical sampling and exchangeability assumptions stated in the underlying theory. Equal engine weighting does not establish these assumptions for ordered, lifetime-stratified degradation trajectories. This study therefore evaluates marginal and subgroup coverage empirically. It does not claim guaranteed coverage for each engine, lifecycle stage, operating regime, complete trajectory, or adaptively selected maintenance time. Nominal 90% coverage is a target, not a guaranteed realized coverage proportion. Maintenance outcomes are evaluated separately under the frozen simulator.

Despite this theoretical appeal, applying conformal prediction to industrial asset prognostics reveals critical research challenges:
\begin{enumerate}
  \item \textbf{Autocorrelated Degradation Trajectories}: Telemetry cycles within an engine violate observation-level exchangeability. Standard pooled conformal prediction over-weights long-lived engines, introducing survivor bias. Hierarchical, engine-weighted quantile computation~\cite{lee2023hierarchical} balances weight across engine clusters, but finite-sample reliability under trajectory autocorrelation requires empirical validation.
  \item \textbf{Subgroup Reliability Breakdown}: Global conformal calibrators target marginal coverage averaged across all operating cycles, but frequently exhibit acute localized undercoverage during high-stress flight regimes or during accelerated terminal wear stages~\cite{yang2026empirical,benabdennour2026grouped}.
  \item \textbf{The Downstream Decision Gap}: Downstream sequential maintenance policies trigger replacement actions when the lower prediction bound drops below an operational threshold ($L(t) \le 1.0$). Conformal calibration methods that treat intervals symmetrically or arbitrarily expand interval width can induce extreme operational conservatism, scrapping components prematurely under the guise of high coverage.
\end{enumerate}

To bridge this gap, this paper investigates \textit{decision-aware, condition-adaptive conformal calibration} for turbofan RUL prognostics. We present an offline, trace-driven prognostic decision-support research prototype framework integrating causal temporal windowing, frozen point-predictor benchmarks, exact rational hierarchical quantiles, and asymmetric nonconformity score inversion, evaluated within a discrete-event sequential maintenance simulator. Physical engine integration and operational airline validation have not been demonstrated, and maintenance costs are evaluated in Simulated Cost Units (SCU) reflecting synthetic benchmark preferences rather than validated airline economics.
